% =============================================================================
% preamble.tex
% Configuracion completa de paquetes y estilos para el libro de Fisicoquimica
% =============================================================================

% -----------------------------------------------------------------------------
% CODIFICACION Y LENGUAJE
% -----------------------------------------------------------------------------
\usepackage[utf8]{inputenc}          % Codificacion de entrada UTF-8
\usepackage[T1]{fontenc}             % Codificacion de fuente T1 (acentos correctos)
\usepackage[spanish,es-nolayout]{babel} % Idioma espanol; es-nolayout evita cambios
                                        % automaticos en listas y teoremas

% Renombrar "Tabla" por "Cuadro" (convencion en ciencias en espanol)
\addto\captionsspanish{%
    \renewcommand{\tablename}{Cuadro}%
    \renewcommand{\listtablename}{\'Indice de cuadros}%
    \renewcommand{\contentsname}{\'Indice de contenidos}%
    \renewcommand{\listfigurename}{\'Indice de figuras}%
    \renewcommand{\figurename}{Figura}%
    \renewcommand{\chaptername}{Cap\'itulo}%
    \renewcommand{\appendixname}{Ap\'endice}%
    \renewcommand{\indexname}{\'Indice alfab\'etico}%
    \renewcommand{\bibname}{Bibliograf\'ia}%
}

% -----------------------------------------------------------------------------
% TIPOGRAFIA: PALATINO
% Newpxtext/newpxmath da Palatino de alta calidad con soporte matematico
% -----------------------------------------------------------------------------
\usepackage{newpxtext}               % Palatino para texto
\usepackage[                         %
    varbb,                           % Usar \mathbb de newpx (no de amssymb)
    varg,                            % Variante de g, y cursivas mas naturales
    amsthm,                          % Integrar con amsthm
]{newpxmath}                         % Palatino para matematicas
\usepackage{microtype}               % Mejora el espaciado tipografico (kerning, etc.)

% -----------------------------------------------------------------------------
% MATEMATICAS Y CIENCIAS
% -----------------------------------------------------------------------------
\usepackage{amsmath}                 % Entornos matematicos avanzados (align, cases...)
% NOTA: amssymb NO se carga explicitamente porque newpxmath ya lo incluye
% internamente. Cargarlo de nuevo provoca el error:
%   "l.261 ...ol{\Bbbk} {\mathord}{AMSb}{"7C}"
% por redefinicion del simbolo \Bbbk del font AMSb.
% Todos los simbolos de amssymb (\mathbb, \mathfrak, etc.) estan disponibles.
\usepackage{amsthm}                  % Entornos de teoremas
\usepackage[version=4]{mhchem}       % Ecuaciones y formulas quimicas: \ce{H2O}
\usepackage{siunitx}                 % Unidades SI: \SI{9.81}{m/s^2}
\DeclareSIUnit\atm{atm}
\DeclareSIUnit\week{semana}
\DeclareSIUnit{\calory}{cal}
\DeclareSIUnit{\year}{yr}
 \DeclareSIUnit\mmHg{mmHg}
 
\sisetup{
    output-decimal-marker = {.},     % Separador decimal: punto
    inter-unit-product = \ensuremath{{}\cdot{}}, % Punto medio entre unidades
    per-mode = symbol,               % Fraccion de unidades con /
    detect-all,                      % Detectar fuente del texto circundante
}
\decimalpoint



% -----------------------------------------------------------------------------
% GRAFICOS Y FIGURAS
% -----------------------------------------------------------------------------
\usepackage{graphics}                % Inclusion basica de graficos
\usepackage[pdftex]{graphicx}                % Inclusion avanzada: [width=], [angle=], etc.
\usepackage{epstopdf}                % Conversion automatica EPS -> PDF
\usepackage{wrapfig}                 % Figuras con texto envolvente (wrapfigure)
\usepackage{float}                   % Control de posicion de flotantes [H]
\usepackage{caption}                 % Personalizar leyendas de figuras y cuadros
\captionsetup{
    font=small,                      % Fuente de caption mas pequena que el texto
    labelfont=bf,                    % "Figura X.Y" en negritas
    labelsep=period,                 % Separador: "Figura 1.1. Titulo"
    skip=6pt,                        % Espacio entre figura y caption
}
\usepackage{subcaption}              % Sub-figuras: (a), (b)...

% Ruta por defecto para buscar figuras
\graphicspath{{figuras/}{../figuras/}}

% -----------------------------------------------------------------------------
% TABLAS
% -----------------------------------------------------------------------------
\usepackage{booktabs}                % Lineas horizontales profesionales: \toprule etc.
\usepackage{tabularx}                % Tablas con columnas de ancho automatico (X)
\usepackage{threeparttable}          % Tablas con notas al pie: \tnote
\usepackage{supertabular}            % Tablas largas que abarcan varias paginas
\usepackage{array}                   % Opciones avanzadas de columnas en tablas
\usepackage{multirow}                % Celdas que abarcan multiples filas

% -----------------------------------------------------------------------------
% DISEÑO DE PAGINA Y ESPACIADO
% -----------------------------------------------------------------------------
\usepackage{geometry}
\geometry{
    paperwidth  = 8in,               % Ancho de pagina: 6 pulgadas
    paperheight = 11in,               % Alto de pagina: 9 pulgadas
    top         = 1in,               % Margen superior
    bottom      = 1in,               % Margen inferior
    inner       = 1in,               % Margen interior (encuadernacion)
    outer       = 0.75in,            % Margen exterior
    includeheadfoot,                 % Incluir encabezado y pie en el calculo
    headheight  = 14pt,              % Altura del encabezado
}
\usepackage{setspace}                % Control de interlineado
\setstretch{1.05}                    % Interlineado ligeramente aumentado (legibilidad)
\usepackage{parskip}                 % Parrafos con espacio vertical en lugar de sangria
\setlength{\parskip}{0.5em}

% -----------------------------------------------------------------------------
% ENCABEZADOS Y PIES DE PAGINA
% -----------------------------------------------------------------------------
\usepackage{fancyhdr}
\pagestyle{fancy}
\fancyhf{}                           % Limpiar todos los campos
\fancyhead[LE]{\small\itshape\leftmark}   % Izquierda paginas pares: capitulo
\fancyhead[RO]{\small\itshape\rightmark}  % Derecha paginas impares: seccion
\fancyfoot[C]{\small\thepage}             % Centro: numero de pagina
\renewcommand{\headrulewidth}{0.4pt}      % Linea bajo el encabezado
\renewcommand{\footrulewidth}{0pt}        % Sin linea en el pie

% Estilo para paginas de inicio de capitulo (plain)
\fancypagestyle{plain}{
    \fancyhf{}
    \fancyfoot[C]{\small\thepage}
    \renewcommand{\headrulewidth}{0pt}
}

% -----------------------------------------------------------------------------
% ESTILOS DE CAPITULO: usando titlesec
% -----------------------------------------------------------------------------
\usepackage{titlesec}
\usepackage{titletoc}

% Formato del titulo de capitulo
\titleformat{\chapter}[display]
    {\normalfont\huge\bfseries\color{chaptercolor}}   % Formato del titulo
    {\chaptertitlename\ \thechapter}                  % Etiqueta "Capitulo N"
    {10pt}                                            % Espacio entre etiqueta y titulo
    {\Huge}                                           % Formato del nombre del capitulo

\titlespacing*{\chapter}{0pt}{-20pt}{30pt}           % Espaciado alrededor del titulo

% Formato de secciones
\titleformat{\section}{\large\bfseries\color{sectioncolor}}{\thesection}{1em}{}
\titleformat{\subsection}{\normalsize\bfseries}{\thesubsection}{1em}{}
\titleformat{\subsubsection}{\normalsize\itshape}{\thesubsubsection}{1em}{}

% -----------------------------------------------------------------------------
% COLORES
% -----------------------------------------------------------------------------
\usepackage{xcolor}
\definecolor{chaptercolor}{RGB}{0, 70, 127}      % Azul marino para titulos de capitulo
\definecolor{sectioncolor}{RGB}{0, 100, 160}     % Azul para secciones
\definecolor{examplecolor}{RGB}{0, 110, 84}      % Verde para ejemplos
\definecolor{exercisecolor}{RGB}{140, 50, 0}     % Cafe para ejercicios
\definecolor{notecolor}{RGB}{180, 100, 0}        % Naranja para notas
\definecolor{linkcolor}{RGB}{0, 70, 127}         % Color de hipervinculos

% -----------------------------------------------------------------------------
% CAJAS Y ENTORNOS ESPECIALES (tcolorbox)
% -----------------------------------------------------------------------------
\usepackage[most]{tcolorbox}
\tcbuselibrary{breakable, skins}

% Entorno: Ejemplo resuelto
\newtcolorbox[auto counter, number within=chapter]{ejemplo}[2][]{%
    enhanced,
    breakable,
    colback  = examplecolor!5,
    colframe = examplecolor,
    fonttitle= \bfseries,
    title    = Ejemplo~\thetcbcounter\ifstrempty{#2}{}{{}: #2},
    label    = ej:\thetcbcounter,
    #1
}

% Entorno: Nota al margen / Nota importante
\newtcolorbox{nota}[1][Nota]{%
    enhanced,
    breakable,
    colback  = notecolor!8,
    colframe = notecolor,
    fonttitle= \bfseries,
    title    = #1,
    left     = 6pt,
    right    = 6pt,
}

% Entorno: Ejercicios (bloque al final de seccion)
\newtcolorbox{bloque_ejercicios}[1][Ejercicios]{%
    enhanced,
    breakable,
    colback  = exercisecolor!5,
    colframe = exercisecolor,
    fonttitle= \bfseries\large,
    title    = #1,
    left     = 8pt,
    right    = 8pt,
    top      = 6pt,
    bottom   = 6pt,
}
% -----------------------------------------------------------------------------
% ENTORNOS DE REACCIONES
%
\makeatletter
\newcounter{reaction}
%%% >> for article <<
%%\renewcommand\thereaction{C\,\arabic{reaction}}
%%% << for article <<
%%% >> for report and book >>
\renewcommand\thereaction{C\,\thechapter.\arabic{reaction}}
\@addtoreset{reaction}{chapter}
%%% << for report and book <<
\newcommand\reactiontag%
  {\refstepcounter{reaction}\tag{\thereaction}}
\newcommand\reaction@[2][]%
  {\begin{equation}\ce{#2}%
  \ifx\@empty#1\@empty\else\label{#1}\fi%
  \reactiontag\end{equation}}
\newcommand\reaction@nonumber[1]%
   {\begin{equation*}\ce{#1}\end{equation*}}
\newcommand\reaction%
  {\@ifstar{\reaction@nonumber}{\reaction@}}
\makeatother
% -----------------------------------------------------------------------------
% ENTORNOS DE TEOREMAS / DEFINICIONES
% -----------------------------------------------------------------------------
\theoremstyle{definition}
\newtheorem{definicion}{Definici\'on}[chapter]
\newtheorem{ley}{Ley}[chapter]
\newtheorem{postulado}{Postulado}[chapter]

\theoremstyle{plain}
\newtheorem{teorema}{Teorema}[chapter]
\newtheorem{corolario}{Corolario}[chapter]

\theoremstyle{remark}
\newtheorem{observacion}{Observaci\'on}[chapter]

% -----------------------------------------------------------------------------
% LISTAS Y ENUMERACIONES
% -----------------------------------------------------------------------------
\usepackage{enumitem}
\setlist[itemize]{noitemsep, topsep=4pt}
\setlist[enumerate]{noitemsep, topsep=4pt}

% Lista de ejercicios numerados
\newlist{ejercicios}{enumerate}{1}
\setlist[ejercicios]{
    label     = \textbf{\arabic*.},
    leftmargin= 2em,
    itemsep   = 6pt,
    topsep    = 4pt,
}

% -----------------------------------------------------------------------------
% GLOSARIO
% -----------------------------------------------------------------------------
\usepackage[
    toc,                             % Agregar glosario al indice de contenidos
    acronym,                         % Seccion separada para acronimos
    style=altlist,                   % Estilo: termino + descripcion en parrafo
    nonumberlist,                    % Sin numeros de pagina en el glosario
]{glossaries}
\newglossary[slg]{symbols}{syi}{syg}{Lista de símbolos}
\makeglossaries
%\setglossarystyle{list}
% -----------------------------------------------------------------------------
% HIPERVINCULOS (hyperref debe ir casi al final del preambulo)
% -----------------------------------------------------------------------------
\usepackage{hyperref}
\hypersetup{
    pdfauthor    = {Dr. José Agustín García Reynoso},
    pdftitle     = {Procesos y Reacciones },
    pdfsubject   = {Atmosf\'ericas},
    pdfkeywords  = {fisicoquímica termodinámica cinética química},
    colorlinks   = true,
    linkcolor    = linkcolor,
    citecolor    = linkcolor,
    urlcolor     = linkcolor,
    bookmarks    = true,
    bookmarksopen= true,
    pdfdisplaydoctitle = true,
    pdfstartview = FitH,
}

\usepackage{bookmark}                % Mejor manejo de bookmarks PDF

% -----------------------------------------------------------------------------
% NUMERACION DE LINEAS (util para revision editorial; comentar en version final)
% -----------------------------------------------------------------------------
\usepackage{lineno}
% Descomentar la siguiente linea para activar numeracion de lineas:
% \linenumbers

% -----------------------------------------------------------------------------
% DIAGRAMAS CON EPIC (entornos picture extendidos)
% -----------------------------------------------------------------------------
\usepackage{epic}                    % Extiende el entorno picture de LaTeX

% -----------------------------------------------------------------------------
% REFERENCIAS Y CITAS
% -----------------------------------------------------------------------------
\usepackage[numbers, sort&compress]{natbib}  % Referencias numeradas y comprimidas
\bibliographystyle{achemso}          % Estilo ACS (quimica); cambiar segun editorial

% -----------------------------------------------------------------------------
% INDICE ANALITICO
% -----------------------------------------------------------------------------
\usepackage{makeidx}
\makeindex

% -----------------------------------------------------------------------------
% UTILIDADES GENERALES
% -----------------------------------------------------------------------------
\usepackage{xspace}                  % Espaciado inteligente despues de macros
\usepackage{ifthen}                  % Condicionales en LaTeX
\usepackage{calc}                    % Calculos aritmeticos en LaTeX
\usepackage{etoolbox}                % Herramientas de programacion LaTeX

% -----------------------------------------------------------------------------
% MACROS DE CONVENIENCIA
% Atajos para simbolos y unidades frecuentes en fisicoquimica
% -----------------------------------------------------------------------------

% Constantes fisicas (valor + unidad en formato siunitx)
\newcommand{\NA}{\ensuremath{N_{\!A}}\xspace}    % Numero de Avogadro
\newcommand{\kb}{\ensuremath{k_B}\xspace}         % Constante de Boltzmann
\newcommand{\Rgas}{\ensuremath{R}\xspace}         % Constante universal de los gases
\newcommand{\hplanck}{\ensuremath{h}\xspace}      % Constante de Planck

% Notacion de derivadas
\newcommand{\dif}{\mathop{}\!\mathrm{d}}          % Diferencial: d (redondo)
\newcommand{\Dif}{\mathop{}\!\mathrm{D}}          % Derivada funcional
\newcommand{\pder}[2]{\frac{\partial #1}{\partial #2}}   % Derivada parcial
\newcommand{\der}[2]{\frac{\dif #1}{\dif #2}}            % Derivada total

% Notacion termodinamica
\newcommand{\DeltaH}{\ensuremath{\Delta H}\xspace}
\newcommand{\DeltaG}{\ensuremath{\Delta G}\xspace}
\newcommand{\DeltaS}{\ensuremath{\Delta S}\xspace}
\newcommand{\Keq}{\ensuremath{K_{\text{eq}}}\xspace}
\newcommand{\stdstate}{\ensuremath{^{\circ}}\xspace}  % Superindice de estado estandar

% Abreviaciones utiles
\newcommand{\eg}{e.g.\xspace}
\newcommand{\ie}{i.e.\xspace}
\newcommand{\etal}{\textit{et al.}\xspace}

%%%%% Problem Set Macros

% =============================================================================
% FIN DE preamble.tex
% =============================================================================
